\section*{Примечания о чтениях}
\subsection*{\texttt{GELFAND-P001485}}
\par\noindent\leavevmode\phantomsection\label{GELFAND-0161-U001}\label{GELFAND-P001485:eq10-cancellation}\textbf{\texttt{GELFAND-0161-U001}} (eq10-cancellation).\nopagebreak[4]\par
В показателе экспоненты (10), после второго интеграла и над левым краем уцелевшего m, находится маленький густо зачёркнутый знак. Уцелевшая запись читается mλ dt; ниже явно написано λ² = ẋ²/m².\par
Кандидаты: зачёркнуто m; зачёркнут иной одиночный знак. Чтение не установлено в дипломатическом исходнике.\par
Координаты исходного изображения: (1805,60,1870,126); более широкая область формулы: (70,35,2210,230); следующее стационарное соотношение: (1950,310,2230,475). Размер кадра 2408×3480. Просмотрены полный оригинал, контекст с увеличением 3× и зачёркнутый знак с увеличением 7×. Штрихи скрывают зачёркнутый символ. Формальное дифференцирование уцелевшего выражения λ⁻¹ẋ² + mλ дало бы знаменатель m в стационарном соотношении, тогда как в источнике написано m². Оба исходных выражения сохранены; m² и числовой коэффициент в (10) не внесены.\par
Повторная сверка: Зачёркнутый знак не восстановлен при 5×/10×; уцелевшая m и отдельный знаменатель m² подтверждены. Изменённый инструментом широкий контекст заменён двумя перекрывающимися точными просмотрами 3×.\par\medskip
\par\noindent\leavevmode\phantomsection\label{GELFAND-0161-U002}\label{GELFAND-P001485:eq11-differential}\textbf{\texttt{GELFAND-0161-U002}} (eq11-differential).\nopagebreak[4]\par
Последний дифференциал в показателе (11) имеет явно x-подобную последнюю букву; в соседней формуле (10) t написана иначе.\par
Кандидаты: dx; dt. Чтение не установлено в дипломатическом исходнике.\par
Координаты исходного изображения: (980,460,1090,540). Размер кадра 2408×3480. Просмотрены полный оригинал, формула с увеличением 3× и дифференциал с увеличением 7×. Визуально сильнее чтение dx; оно и воспроизведено. Привычное выражение с временным параметром подсказывало бы dt, но им не заменена видимая запись.\par
Повторная сверка: Форма x по-прежнему предпочтительнее при 5×/10×; вариант dt сохранён без исправления формулы.\par\medskip
\par\noindent\leavevmode\phantomsection\label{GELFAND-0161-U003}\label{GELFAND-P001485:metric-abbreviation}\textbf{\texttt{GELFAND-0161-U003}} (metric-abbreviation).\nopagebreak[4]\par
В скобках после «1-мерн. метрики» читается «всё остальное … из репараметризации». Сокращение перед «из» выглядит как «свод.» либо «своб.»; последняя буква имеет петлю.\par
Кандидаты: свод.; своб.. Чтение не установлено в дипломатическом исходнике.\par
Координаты исходного изображения: (1970,764,2120,860); более широкая область скобок: (1550,755,2270,930). Размер кадра 2408×3480. Просмотрены контекст с увеличением 3×, слово с увеличением 7× и слово вместе со следующим «из» с увеличением 5×. Последняя рукописная буква похожа на д, но её нельзя уверенно отличить от б этого почерка; сокращение не раскрыто.\par
Повторная сверка: Петля скорее допускает д, но не исключает б; сокращение не раскрыто.\par\medskip
\par\noindent\leavevmode\phantomsection\label{GELFAND-0161-U004}\label{GELFAND-P001485:eq12-identity}\textbf{\texttt{GELFAND-0161-U004}} (eq12-identity).\nopagebreak[4]\par
Формула (12) явно записана как ∫(dλ/λ)e\textasciicircum{}\{aλ\}e\textasciicircum{}\{b/λ\}=e\textasciicircum{}\{√(ab)\}. Пределы интегрирования, дополнительные множители и коэффициент 2 перед корнем не видны.\par
Кандидаты: \textbackslash{}int\textbackslash{}frac\{d\textbackslash{}lambda\}\{\textbackslash{}lambda\}e\textasciicircum{}\{a\textbackslash{}lambda\}e\textasciicircum{}\{b/\textbackslash{}lambda\}=e\textasciicircum{}\{\textbackslash{}sqrt\{ab\}\}. Чтение не установлено в дипломатическом исходнике.\par
Координаты исходного изображения: (290,864,1200,1070). Размер кадра 2408×3480. Просмотрены оригинал и фрагмент с увеличением 3×. Равенство сохранено как написано. Его аналитический смысл, условия сходимости и нормировка источником не установлены; математически ожидаемый множитель или другая функция не внесены.\par
Повторная сверка: Первоначальные детали обрезали низ меры и верх корня. Исправленные 5×/10× показывают весь знаменатель λ и черту корня; условный сходящийся интеграл вычислен отдельно.\par\medskip
\par\noindent\leavevmode\phantomsection\label{GELFAND-0161-U005}\label{GELFAND-P001485:eq14-differentials}\textbf{\texttt{GELFAND-0161-U005}} (eq14-differentials).\nopagebreak[4]\par
В (14) после высоко написанного ∫G\_\{μν\}(x)√g g\textasciicircum{}\{ab\}∂\_a x\textasciicircum{}μ∂\_b x\textasciicircum{}ν буквы dσ dt помещены справа на более низкой строке, примерно на уровне основной экспоненты e.\par
Кандидаты: d\textbackslash{}sigma\textbackslash{},dt принадлежит интегралу в показателе экспоненты; d\textbackslash{}sigma\textbackslash{},dt написано вне показателя экспоненты. Чтение не установлено в дипломатическом исходнике.\par
Координаты исходного изображения: (850,1790,1840,2020); дифференциалы: (1640,1920,1840,2020). Размер кадра 2408×3480. Просмотрены оригинал и фрагмент с увеличением 3×. В транскрипции сохранено видимое низкое расположение во второй строке выровненной формулы. Математический контекст скорее указывает на принадлежность дифференциалов интегралу, но одна исходная компоновка не позволяет окончательно установить область действия показателя.\par
Повторная сверка: Низкое расположение dσ dt подтверждено; область действия показателя не восстановлена из математического ожидания.\par\medskip
\subsection*{\texttt{GELFAND-P001486}}
\par\noindent\leavevmode\phantomsection\label{GELFAND-0161-U006}\label{GELFAND-P001486:feynman-verb}\textbf{\texttt{GELFAND-0161-U006}} (feynman-verb).\nopagebreak[4]\par
В реплике И.М.Г. между «и» и словом «эстетике» видна короткая слитная запись «учи…»; её окончание сжато.\par
Кандидаты: учить; учите; учиться. Чтение не установлено в дипломатическом исходнике.\par
Координаты исходного изображения: (950,2110,1145,2200). Размер кадра 2408×3480. Просмотрены полный оригинал, строка с увеличением 3× и фрагмент с увеличением 7×. Визуально предпочтительно «учить»; «учиться» — лишь возможное чтение сжатых штрихов окончания и не принято по грамматическим соображениям.\par
Повторная сверка: Сжатое окончание не разрешено при 5×/10×; предпочтительный исторический вариант не объявлен установленным.\par\medskip
\par\noindent\leavevmode\phantomsection\label{GELFAND-0161-U007}\label{GELFAND-P001486:feynman-noun}\textbf{\texttt{GELFAND-0161-U007}} (feynman-noun).\nopagebreak[4]\par
Следующее слово в реплике о Фейнмане выглядит как «эстетике»; последний соединённый штрих допускает чтение «эстетику». После него отчётливо «у него».\par
Кандидаты: эстетике; эстетику. Чтение не установлено в дипломатическом исходнике.\par
Координаты исходного изображения: (1125,2110,1330,2200). Размер кадра 2408×3480. Просмотрены полный оригинал, строка с увеличением 3× и фрагмент с увеличением 7×. Форма последней буквы, похожая на е, визуально убедительнее, но сжатое окончание отмечено отдельно от предыдущего глагола.\par
Повторная сверка: Окончание е предпочтительнее, но у не исключено; следующие «у него» читаются независимо.\par\medskip
\par\noindent\leavevmode\phantomsection\label{GELFAND-0161-U008}\label{GELFAND-P001486:names-cancellation}\textbf{\texttt{GELFAND-0161-U008}} (names-cancellation).\nopagebreak[4]\par
После «Поляков и» перед «Замолодчиков — Книжник — Поляков» густо зачёркнуты примерно две буквы. Под горизонтальными штрихами просматривается подобие «по».\par
Кандидаты: по; По. Чтение не установлено в дипломатическом исходнике.\par
Координаты исходного изображения: (635,2568,755,2645). Размер кадра 2408×3480. Просмотрены оригинал, строка имён с увеличением 3× и зачёркнутое место с увеличением 10×. Закрывающие штрихи не позволяют надёжно восстановить зачёркнутый текст и регистр букв. Незачёркнутые имена источника и их раздельное подчёркивание сохранены.\par
Повторная сверка: Оставшиеся петли не устанавливают регистр. Порядок имён источника и отдельно записанный английский алфавитный порядок сохранены.\par\medskip
\par\noindent\leavevmode\phantomsection\label{GELFAND-0161-U009}\label{GELFAND-P001486:metric-diagram}\textbf{\texttt{GELFAND-0161-U009}} (metric-diagram).\nopagebreak[4]\par
Внизу справа видны пересекающиеся горизонтальная и вертикальная оси с подписями ξ₁ и ξ₂. В правой верхней области расположен неровный замкнутый контур с g\_ab(ξ) внутри. Утолщённые и заострённые концы осей не позволяют уверенно восстановить все направления стрелок.\par
Кандидаты: Горизонтальная ось со стрелками влево и вправо.; Горизонтальная ось со стрелкой вправо; левое утолщение — след пера.; Утолщения на концах осей не задают уверенно читаемых направлений.. Чтение не установлено в дипломатическом исходнике.\par
Координаты исходного изображения: (1650,3105,2170,3420); левый и вертикальный концы осей: (1725,3160,1830,3390); правый конец: (2010,3270,2120,3350). Размер кадра 2408×3480. Просмотрены полный оригинал, рисунок с увеличением 3×, левый и вертикальный концы с увеличением 8× и правый конец с увеличением 10×. Все подписи читаются. Маркер SourceFigure сохраняет их расположение, пересечение осей и подпись метрики внутри контура; неясные наконечники оставлены для проверки по источнику.\par
Повторная сверка: Три подписи ясны и набраны. Малые концевые крючки оставлены штрихами без приписанного направления; рядом сохранён оригинал.\par\medskip
\subsection*{\texttt{GELFAND-P001487}}
\par\noindent\leavevmode\phantomsection\label{GELFAND-0161-U101}\label{GELFAND-P001487:top-beta-inequality}\textbf{\texttt{GELFAND-0161-U101}} (top-beta-inequality).\nopagebreak[4]\par
В строке «Интеграл (2) опр. при ...» между β и β\_c густо переправленный знак: наружные штрихи образуют >, внутри видны дополнительные штрихи.\par
Кандидаты: β > β\_c; β ≥ β\_c. Чтение не установлено в дипломатическом исходнике.\par
Координаты оригинала 2408×3480: (1075,62,1270,165). Проверено на полном оригинале, при увеличении 3× и отдельном увеличении 10×. Возможное предшествующее чтение затёрто; варианты не утверждаются.\par
Повторная сверка: Наружный > и внутренние штрихи видны, но строгий или нестрогий знак не установлен.\par\medskip
\par\noindent\leavevmode\phantomsection\label{GELFAND-0161-U102}\label{GELFAND-P001487:left-margin-monogram}\textbf{\texttt{GELFAND-0161-U102}} (left-margin-monogram).\nopagebreak[4]\par
На левом поле напротив «Всё совпало!!» помещена короткая крупная рукописная помета, похожая на слитные инициалы.\par
Кандидаты: ИГ; ИМ. Чтение не установлено в дипломатическом исходнике.\par
Координаты оригинала 2408×3480: (120,1200,248,1372). Полный оригинал, 3× и 10×; буквы и возможность небуквенной пометы окончательно не различены.\par
Повторная сверка: Монограмма и небуквенная помета остаются возможными; атрибуция не выведена.\par\medskip
\par\noindent\leavevmode\phantomsection\label{GELFAND-0161-U103}\label{GELFAND-P001487:upper-dimension-range}\textbf{\texttt{GELFAND-0161-U103}} (upper-dimension-range).\nopagebreak[4]\par
После формулы (4) отчётливо написано «Это имеет смысл при D ≤ 1 и D ≥ 26». В самой формуле стоит корень из (D−1)(D−25).\par
Кандидаты: D ≤ 1 и D ≥ 26 — буквальная запись источника. Чтение не установлено в дипломатическом исходнике.\par
Координаты оригинала 2408×3480: (330,715,1730,810); формула для сравнения (350,560,1750,710). Проверено на оригинале и в увеличении 3×. Алгебраическая проверка подкоренного выражения даёт также границу D=25; число 26 в рукописи оставлено без исправления. Причина выбора 26 в источнике не установлена.\par
Повторная сверка: Подтверждены 26 в фразе, 25 под корнем и союз «и». Область вещественности вычислена отдельно.\par\medskip
\par\noindent\leavevmode\phantomsection\label{GELFAND-0161-U104}\label{GELFAND-P001487:middle-cancelled-word}\textbf{\texttt{GELFAND-0161-U104}} (middle-cancelled-word).\nopagebreak[4]\par
Перед словом «Асимптотика» находится слово, густо зачёркнутое двумя наклонными линиями; нижние части букв похожи на «число».\par
Кандидаты: число. Чтение не установлено в дипломатическом исходнике.\par
Координаты оригинала 2408×3480: (356,2195,571,2266). Проверено при 3× и 10×; часть букв пересечена густыми чертами, поэтому чтение оставлено вопросительным.\par
Повторная сверка: Начальные и конечные петли не дают полного отменённого слова.\par\medskip
\par\noindent\leavevmode\phantomsection\label{GELFAND-0161-U105}\label{GELFAND-P001487:right-margin-triangle-graphs}\textbf{\texttt{GELFAND-0161-U105}} (right-margin-triangle-graphs).\nopagebreak[4]\par
Два треугольных рисунка на правом поле. Верхний: внешний треугольник, внутренний треугольник и три соединения соответствующих вершин. На нижнем внутреннем ребре слева и около правой внешней вершины есть утолщённые направленные штрихи. Нижний: внешний треугольник с тремя отрезками к внутренней точке; левый нижний отрезок заканчивается с небольшим разрывом перед этой точкой. Подписей нет.\par
Кандидаты: Утолщённые штрихи верхнего рисунка обозначают стрелки: на внутреннем основании влево, у внешнего правого угла вправо.; Утолщённые штрихи — следы повторного проведения рёбер, не стрелки.; Разрыв в нижнем рисунке — разрыв пера при изображении общей внутренней вершины.; Разрыв в нижнем рисунке намеренно различает окончания рёбер.. Чтение не установлено в дипломатическом исходнике.\par
Координаты оригинала 2408×3480: (1740,1740,2140,2300). Проверены полный оригинал и 3×; выполнен также независимый повторный просмотр увеличений. Чёткая геометрия описана в SourceFigure; направление и значение сомнительных штрихов не дорисованы.\par
Повторная сверка: Наконечники и повторная обводка не различены; нижний центральный разрыв ясен. Воссоздание сохраняет разрывы и крючки без добавления узла.\par\medskip
\par\noindent\leavevmode\phantomsection\label{GELFAND-0161-U106}\label{GELFAND-P001487:bottom-cancelled-fraction}\textbf{\texttt{GELFAND-0161-U106}} (bottom-cancelled-fraction).\nopagebreak[4]\par
Перед оставленной правой частью (8) зачёркнута дробь с левым множителем, похожим на 2. В числителе читается (4n+7)! либо близкая запись; n/k, плюс/минус и 7/1 пересечены штрихами. В знаменателе видны скобки и запись, похожая на n−4; правый факториал неотличим от зачёркивания.\par
Кандидаты: 2(4n+7)!/(n−4); 2(4n+1)!/(n−4); 2(4n−7)!/(n−4); 2(4k+7)!/(k−4)!. Чтение не установлено в дипломатическом исходнике.\par
Координаты оригинала 2408×3480: (544,2840,967,3060). Проверено на полном оригинале, при 3× и 8×; при независимом повторном просмотре проверена область (535,2830,975,3070) при 8×. В транскрипции сохранены сама отмена, дробная структура и варианты всех неясных позиций. Эта отменённая дробь не заменена сохранённой справа формулой. Ни один полный вариант не утверждается.\par
Повторная сверка: Первоначальные детали обрезали части дроби. Исправленный общий просмотр 5× и перекрывающиеся детали 10× покрывают числитель и знаменатель; полный вариант всё ещё не установлен.\par\medskip
\par\noindent\leavevmode\phantomsection\label{GELFAND-0161-U107}\label{GELFAND-P001487:bottom-final-factor}\textbf{\texttt{GELFAND-0161-U107}} (bottom-final-factor).\nopagebreak[4]\par
Последний множитель в оставленной формуле (8) имеет верхний индекс k. Его основание — петля с длинным загнутым вниз штрихом, похожая как на 9, так и на латинское g.\par
Кандидаты: 9\textasciicircum{}k; g\textasciicircum{}k. Чтение не установлено в дипломатическом исходнике.\par
Координаты оригинала 2408×3480: (1390,2870,1540,3012). Проверены оригинал, 3× и 10×; независимый повторный просмотр также оставил оба чтения неразрешёнными. Буквальная версия 9\textasciicircum{}k дала бы дополнительный геометрический множитель по сравнению с (7), но это математическое соображение не разрешает графическую неоднозначность и не использовано для исправления источника.\par
Повторная сверка: Полная петля и нижний штрих допускают 9 или g. Асимптотика не устанавливает написанное основание.\par\medskip
\subsection*{\texttt{GELFAND-P001488}}
\par\noindent\leavevmode\phantomsection\label{GELFAND-0161-U110}\label{GELFAND-P001488:top-triangle-graph}\textbf{\texttt{GELFAND-0161-U110}} (top-triangle-graph).\nopagebreak[4]\par
Слева вверху внешний треугольник и внутренний треугольник с подписями 1 (верхняя внутренняя вершина), i (левая нижняя), 2 (правая нижняя). Соответствующие внешние и внутренние вершины соединены. Ребро 1–2 изогнуто. На левой внешней стороне имеется утолщение с коротким боковым штрихом, похожее на направленный вниз наконечник.\par
Кандидаты: На левой внешней стороне обозначена стрелка вниз к левой нижней вершине.; Утолщение на левой внешней стороне — наложение штрихов при повторном проведении ребра.. Чтение не установлено в дипломатическом исходнике.\par
Координаты оригинала 2408×3480: (305,185,680,500). Проверены полный оригинал и 3×; подробный фрагмент также независимо проверен повторно. Чёткие подписи и геометрия перечислены в SourceFigure, сомнительный наконечник не нормализован.\par
Повторная сверка: Три подписи набраны, ребро 1–2 изогнуто. Сохранены видимые малые разрывы; наружному крючку направление не приписано.\par\medskip
\par\noindent\leavevmode\phantomsection\label{GELFAND-0161-U111}\label{GELFAND-P001488:middle-cancelled-incidence}\textbf{\texttt{GELFAND-0161-U111}} (middle-cancelled-incidence).\nopagebreak[4]\par
В подписи к G после «Матр.» зачёркнуто длинное слово, похожее на «инцидентности»; ниже оставлено «смежности вершин».\par
Кандидаты: инцидентности. Чтение не установлено в дипломатическом исходнике.\par
Координаты оригинала 2408×3480: (1555,1290,2010,1390). Проверено на полном оригинале, при 3× и 6×. Часть контуров слова скрыта пересекающими его линиями; исправление источника и сомнение в отменённом слове сохранены.\par
Повторная сверка: Длинное зачёркивание скрывает часть букв; действующая подпись о смежности ясна.\par\medskip
\par\noindent\leavevmode\phantomsection\label{GELFAND-0161-U112}\label{GELFAND-P001488:middle-overdrawn-matrix-symbol}\textbf{\texttt{GELFAND-0161-U112}} (middle-overdrawn-matrix-symbol).\nopagebreak[4]\par
В первой строке определения матрицы знак перед индексом ij многократно обведён: окончательная внешняя форма читается как G, но внутри видны дополнительные дуги, похожие на E либо прежнее G.\par
Кандидаты: Прежний знак E, исправленный на G; Повторно обведённый знак G. Чтение не установлено в дипломатическом исходнике.\par
Координаты оригинала 2408×3480: (1305,1440,1610,1560). Проверены оригинал, 3× и 6×. В основном тексте сохранено G\_ij=1; отдельная помета учитывает неопределённость предшествующих штрихов.\par
Повторная сверка: Окончательная G ясна при 5×/10×; прежний знак не восстановлен.\par\medskip
\par\noindent\leavevmode\phantomsection\label{GELFAND-0161-U113}\label{GELFAND-P001488:middle-overdrawn-equivalence}\textbf{\texttt{GELFAND-0161-U113}} (middle-overdrawn-equivalence).\nopagebreak[4]\par
После G\_ij=0 стоит двусторонняя двойная стрелка, пересечённая несколькими густыми наклонными штрихами. Неясно, знак усиленно обведён или отменён.\par
Кандидаты: Действующий знак ⇔, несколько раз обведённый; Зачёркнутый знак ⇔. Чтение не установлено в дипломатическом исходнике.\par
Координаты оригинала 2408×3480: (1560,1560,1715,1650). Проверены полный оригинал, 3× и 10×. Геометрическая форма стрелки видима, но редакционный смысл густых штрихов не установлен.\par
Повторная сверка: Геометрия стрелки видна; её действующий или отменённый статус не установлен.\par\medskip
\par\noindent\leavevmode\phantomsection\label{GELFAND-0161-U114}\label{GELFAND-P001488:lower-cancelled-space-token}\textbf{\texttt{GELFAND-0161-U114}} (lower-cancelled-space-token).\nopagebreak[4]\par
Между словами «векторное» и «пространство» зачёркнут короткий фрагмент; под штрихами видны части двух петель либо букв, возможно начало «пр».\par
Кандидаты: пр; Короткий неразборчивый буквенный или математический знак. Чтение не установлено в дипломатическом исходнике.\par
Координаты оригинала 2408×3480: (680,1968,812,2052). Проверены оригинал, 3× и 10×. Нельзя установить, было ли это начало повторённого слова или отдельный символ; отмена и кандидаты перенесены явно.\par
Повторная сверка: Плотные косые линии скрывают середину; фрагмент не раскрыт до слова.\par\medskip
\par\noindent\leavevmode\phantomsection\label{GELFAND-0161-U115}\label{GELFAND-P001488:lower-tautological-annotation}\textbf{\texttt{GELFAND-0161-U115}} (lower-tautological-annotation).\nopagebreak[4]\par
Под левым O(−1) в правой приписке нарисованы два почти параллельных вертикальных дугообразных штриха, открытых вправо; справа на их уровне стоит \{(ℓ,x)|x∈ℓ\}.\par
Кандидаты: Вертикально поставленный знак равенства, связывающий O(−1) с указанным множеством; Двойная дугообразная помета без однозначно читаемого математического знака. Чтение не установлено в дипломатическом исходнике.\par
Координаты оригинала 2408×3480: (1635,2040,2125,2195). Проверены полный оригинал, 3× и 5×. Формула включения и множество прочитаны уверенно; функция соединяющей их пометы описана без придумывания дополнительного отношения.\par
Повторная сверка: Два дугообразных штриха различимы. Множество и включение набраны без добавленного равенства.\par\medskip
\par\noindent\leavevmode\phantomsection\label{GELFAND-0161-U116}\label{GELFAND-P001488:bottom-cancelled-exercise-symbol}\textbf{\texttt{GELFAND-0161-U116}} (bottom-cancelled-exercise-symbol).\nopagebreak[4]\par
В последней строке между «для» и «dim V=1 и 2» стоит густо зачёркнутый короткий знак, в котором различимы высокие наклонные штрихи.\par
Кандидаты: M; V. Чтение не установлено в дипломатическом исходнике.\par
Координаты оригинала 2408×3480: (1650,3225,1760,3330). Проверены полный оригинал, 3× и 10×. Кандидаты отражают сохранившиеся штрихи и не утверждаются; знак отмены сохранён.\par
Повторная сверка: Высокие штрихи пересечены зачёркиванием; M/V не разрешено.\par\medskip
\subsection*{\texttt{GELFAND-P001489}}
\par\noindent\leavevmode\phantomsection\label{GELFAND-0161-U201}\label{GELFAND-P001489:cancel-after-image}\textbf{\texttt{GELFAND-0161-U201}} (cancel-after-image).\nopagebreak[4]\par
Справа от «Ker v ⊃ Im v,» виден короткий фрагмент, густо зачёркнутый несколькими горизонтальными и петлеобразными штрихами. Исходные знаки не удаётся надёжно отделить от зачёркивания.\par
Кандидаты: не записаны. Чтение не установлено в дипломатическом исходнике.\par
Оригинал 2408×3480; координаты (x0,y0,x1,y1)=(1235,495,1400,566). Этот фрагмент проверен по оригиналу, общему увеличению 3× и отдельному увеличению 10×. Текст под зачёркиванием не восстановлен; ни число букв, ни конкретное слово не утверждаются. В теле сохранено место отменённого фрагмента и маркер Incomplete.\par
Повторная сверка: Полный знак или слово не восстановлены. Маркер Incomplete сохранён отдельно от обычных вариантов чтения.\par\medskip
\par\noindent\leavevmode\phantomsection\label{GELFAND-0161-U202}\label{GELFAND-P001489:cancel-before-bundle}\textbf{\texttt{GELFAND-0161-U202}} (cancel-before-bundle).\nopagebreak[4]\par
Между «Рассм.» и «расслоение» имеется зачёркнутое короткое слово или математическое обозначение. Под плотными линиями различимы отдельные петли, но достоверно прочесть целый фрагмент нельзя.\par
Кандидаты: не записаны. Чтение не установлено в дипломатическом исходнике.\par
Оригинал 2408×3480; координаты (x0,y0,x1,y1)=(527,590,686,679). Выполнены просмотр оригинала, увеличение окружающего текста 3× и отдельное увеличение 10×. Возможное латинское начало не даёт достаточного основания записать слово; место пропуска явно помечено Incomplete.\par
Повторная сверка: Возможного начала по-прежнему недостаточно; слово и число букв не добавлены.\par\medskip
\par\noindent\leavevmode\phantomsection\label{GELFAND-0161-U203}\label{GELFAND-P001489:truncated-arrow-cancel}\textbf{\texttt{GELFAND-0161-U203}} (truncated-arrow-cancel).\nopagebreak[4]\par
Под стрелкой от Λ\textasciicircum{}k V к Λ\textasciicircum{}\{k+1\} V находится небольшой знак, пересечённый наклонным штрихом. Над стрелкой отдельно и ясно написано v.\par
Кандидаты: v; k. Чтение не установлено в дипломатическом исходнике.\par
Оригинал 2408×3480; проверенный фрагмент (x0,y0,x1,y1)=(665,1675,845,1820), увеличение 8×; сам перечёркнутый знак расположен примерно в (703,1755,759,1814). Смысл зачёркнутого знака и его первоначальное чтение не установлены. Надстрочное v и направление стрелки прочитаны независимо; оба кандидата относятся только к отменённому нижнему знаку.\par
Повторная сверка: Нижний знак не разрешён; ясная верхняя v не использована для его восстановления.\par\medskip
\par\noindent\leavevmode\phantomsection\label{GELFAND-0161-U204}\label{GELFAND-P001489:symmetry-note}\textbf{\texttt{GELFAND-0161-U204}} (symmetry-note).\nopagebreak[4]\par
В примере 3 V названо «симплектич. пр-во»; над словами «2-форма» добавлено «симм.» и проведена фигурная соединительная скобка.\par
Кандидаты: симм. (видимое сокращение). Чтение не установлено в дипломатическом исходнике.\par
Оригинал 2408×3480; координаты (x0,y0,x1,y1)=(1785,2135,2110,2265), увеличение 6×, также просмотр всего примера 3×. Чтение «симм.» поддерживается штрихами; перед ним не видно «косо». Соседнее «симплектич.» создаёт возможное математическое несоответствие. Никакое слово не дописано, свойство формы не исправлено; требуется историческая проверка намерения автора.\par
Повторная сверка: Добавленное «симм.» видно без предшествующего «косо». Симплектическая редукция проверена отдельно при явно кососимметричной форме.\par\medskip
\subsection*{\texttt{GELFAND-P001490}}
\par\noindent\leavevmode\phantomsection\label{GELFAND-0161-U205}\label{GELFAND-P001490:basis-abbreviation}\textbf{\texttt{GELFAND-0161-U205}} (basis-abbreviation).\nopagebreak[4]\par
В строке перед e\_1,…,e\_r написано «В орто… . базисе». Конец сокращения имеет плотные соединённые штрихи.\par
Кандидаты: ортон.; ортог.. Чтение не установлено в дипломатическом исходнике.\par
Оригинал 2408×3480; координаты (x0,y0,x1,y1)=(394,1545,605,1664), увеличение 6×; также просмотр общего фрагмента 3×. По форме последней буквы предпочтительно «ортон.», но чтение «ортог.» полностью не исключено. Сокращение не раскрыто и нормировка формы не добавлена.\par
Повторная сверка: Конечная n-подобная форма скорее допускает «ортон.»; полное слово и нормировка не добавлены.\par\medskip
\par\noindent\leavevmode\phantomsection\label{GELFAND-0161-U206}\label{GELFAND-P001490:clifford-normalization}\textbf{\texttt{GELFAND-0161-U206}} (clifford-normalization).\nopagebreak[4]\par
В общем соотношении ясно записано vv′+v′v=⟨v,v′⟩h. Ниже столь же ясно написано e\_i²=h, e\_i e\_k=−e\_k e\_i при i≠k.\par
Кандидаты: vv′+v′v=⟨v,v′⟩h; e\_i²=h (оба соотношения как написаны). Чтение не установлено в дипломатическом исходнике.\par
Оригинал 2408×3480; координаты двух соотношений (x0,y0,x1,y1)=(425,1460,1325,1560) и (384,1662,1399,1777); проверены по оригиналу и в увеличении 3× общего фрагмента (255,1100,2160,1790). Множитель 2 нигде не добавлен и правые части не заменены. Нормировка формы в данном месте не указана; при нормировке ⟨e\_i,e\_i⟩=2 соотношения согласованы. Эта возможная конвенция не утверждается как дополнительный текст источника.\par
Повторная сверка: Первоначальные тесные детали обрезали верхнее соотношение. Исправленная полная полоса 5×, перекрывающиеся детали 10×, прежние нижние детали и целый оригинал подтверждают обе записи. Сопоставление нормировок остаётся условным.\par\medskip
